\input{../tex-inputs/latex-leseansicht-vorspann}

\section[Arthur Schnitzler an Gustav Schwarzkopf, 13. 6. 1928]{L04479 Arthur Schnitzler an Gustav Schwarzkopf, 13. 6. 1928}
\nopagebreak\mylabel{L04479v}
\rehead{ }\normalsize\beginnumbering\briefempfaengerindex{Schwarzkopf, Gustav@\textsc{Schwarzkopf, Gustav}!zzzSchnitzler, Arthur@\emph{von Arthur Schnitzler}!1928-06-131@{13. 6. 1928}|(be}
\toendnotes[C]{\smallbreak\pagebreak[2]}
\correspDesc{Versand  durch Arthur Schnitzler am 13. 6. 1928 in Bad Ischl
\newline{}Übermittlung  am 13. 6. 1928 in Bad Ischl
\newline{}Erhalt  durch Gustav Schwarzkopf im Zeitraum [14. 6. 1928 – 18. 6. 1928?] in Wien}\toendnotes[C]{\smallbreak}
\Standort{DLA, A:Schnitzler, HS.1985.1.1897.}
\physDesc{Bildpostkarte, 353 Zeichen
\newline{}Handschrift: Bleistift, lateinische Kurrent
\newline{}Versand: Stempel: »\nobreak{}\oindex{Bad Ischl@\textbf{Bad Ischl}|pwk}Bad Ischl 1, {[}13{]} \textcolor{gray}{VI.} 28, 19\nobreak{}«.  }\pstart{}{\pb}Hrn Guſtav Schwarzkopf\pend{}\pstart{}Wien I\oindex{I., Innere Stadt@\textbf{I., Innere Stadt}|pw}\pend{}\pstart{}Tiefer Graben 17\oindex{Wien@\textbf{Wien}!I., Innere Stadt@\textbf{I., Innere Stadt}!Tiefer Graben 17@\textbf{Tiefer Graben 17}|pw}.\pend{}{\bigskip}
\pstart
           \noindent{}{\pb}\textcolor{gray}{\textbf{2. Bad Ischl\oindex{Bad Ischl@\textbf{Bad Ischl}|pw} mit Dachstein\oindex{Dachstein@\textbf{Dachstein}|pw}.}}\pend
           \vspace{1em}
\pstart
           {\pb}Ischl\oindex{Bad Ischl@\textbf{Bad Ischl}|pw}. 13. 6. 928\pend
           \vspace{0.5em}
\pstart
           lieber Guſtav, es wäre sehr schön we{\geminationn}
               Sie herkämen. Wohne mit Frau P.\pwindex{Pollaczek, Clara Katharina 15.\,1.\,1875 Wien – 22.\,7.\,1951 ebd.@\textsc{Pollaczek, Clara Katharina} (15.\,1.\,1875 Wien – 22.\,7.\,1951 ebd.)|pw} im Elisabeth\oindex{Hotel Kaiserin Elisabeth [Bad Ischl]@\textbf{Hotel Kaiserin Elisabeth [Bad Ischl]}|pw}  (ausgezeichnet u nicht theuer) – Landschaft
               unverändert schön, Wetter köstlich – und noch wenig Leute. Denke bis 20.
               hier zu bleiben. Also  entschließen Sie sich – es würde {\pb}Ihnen sicher gut thun.\pend
           
\pstart
           Alles herzliche{\\[\baselineskip]}Ihr \spacefill\mbox{A.}\pend
           \leftskip=0em{}\selectlanguage{ngerman}\endnumbering\briefempfaengerindex{Schwarzkopf, Gustav@\textsc{Schwarzkopf, Gustav}!zzzSchnitzler, Arthur@\emph{von Arthur Schnitzler}!1928-06-131@{13. 6. 1928}|)be}\mylabel{L04479h}
\begin{anhang}
\end{anhang}\newcommand{\dateiname}{L04479}\newcommand{\titel}{Arthur Schnitzler an Gustav Schwarzkopf, 13. 6. 1928}\newcommand{\editorInnen}{Herausgegeben von Jahnke, SelmaMüller, Martin Anton}\input{../tex-inputs/latex-leseansicht-abspann}
